%==============================================================================
% Chapter 2 — Theoretical Background & Literature Review
%==============================================================================
\chapter[Theoretical Background and Literature Review]{الخلفية النظرية ومراجعة الأدبيات}
\label{ch:literature}
\markboth{الفصل الثاني: الخلفية النظرية ومراجعة الأدبيات}{الفصل الثاني: الخلفية النظرية ومراجعة الأدبيات}

%------------------------------------------------------------------------------
\section{مقدمة الفصل}
\label{sec:lit-intro}

يستعرض هذا الفصل الأسس النظرية والمراجعة النقدية للبيئة البحثية والتقنية التي
يستند إليها مشروع «المساعد الذكي المؤسسي (بيان)». تتوزع مراجعة الأدبيات في
هذا الفصل على خمسة محاور رئيسية تشمل: النماذج اللغوية الكبيرة (\EN{LLMs})،
وتقنيات تحويل اللغة الطبيعية إلى استعلامات قواعد البيانات (\EN{NL2SQL})،
ونظم الاسترجاع المعزز بالتوليد (\EN{RAG})، والمعماريات القائمة على الأنظمة
متعددة الوكلاء (\EN{Multi-Agent Systems})، بالإضافة إلى تحديات معالجة اللغة
العربية (\EN{Arabic NLP}). ويُختتم الفصل بتحليل مقارن يبرز الفجوة البحثية
والتطبيقية التي يستهدف المشروع معالجتها.

تكتسب هذه المراجعة أهمية مضاعفة في مشروع بيان، نظراً لأنه يدمج أربع تقنيات
متقدمة في إطار موحد. لذا، يتطلب المشروع فحصاً دقيقاً لأسس كل تقنية على حدة،
ثم تحليل كيفية تكاملها في الأدبيات السابقة. كما يسهم هذا الفصل في تحديد الفجوة
البحثية بدقة، وتبرير القرارات التصميمية المُتخذة في الفصل الأول، خاصة اختيار
\EN{LangGraph} لإطار التنسيق بين الوكلاء، واعتماد طبقة \EN{Views} الذكية
لتحسين دقة \EN{NL2SQL}، واستخدام \EN{Odoo ERP} كبيئة محاكاة واقعية.

يتضمن الفصل ستة أقسام رئيسية: يبدأ بالخلفية النظرية للتقنيات الأساسية، ثم
يستعرض الأنظمة والدراسات السابقة، يلي ذلك قسم مقارن يقدم جدولاً تحليلياً، ثم
تحليل الفجوة البحثية، وأخيراً الخلاصة. ويعتمد الفصل على نظام استشهاد \EN{IEEE}
مع استمرارية الترقيم من الفصل الأول.

%------------------------------------------------------------------------------
\section{الخلفية النظرية}
\label{sec:theory}

تستند الخلفية النظرية لمشروع بيان إلى خمس ركائز علمية مترابطة، تشكل مجتمعةً
الأساس النظري للنظام. وفيما يلي عرض تفصيلي لكل ركيزة مع التركيز على التطورات
الحديثة والأدوات التجارية المتاحة.

\subsection{النماذج اللغوية الكبيرة (\EN{Large Language Models - LLMs})}
\label{subsec:llms}

تُعرَّف النماذج اللغوية الكبيرة بأنها نماذج ذكاء اصطناعي ضخمة تُدرَّب على
كميات هائلة من النصوص بهدف فهم اللغة البشرية وتوليدها. وقد شهدت هذه النماذج
تطوراً مذهلاً منذ عام 2020، حيث انتقلت من نماذج متخصصة في مهمة واحدة إلى
نماذج عامة قادرة على أداء آلاف المهام بدقة عالية.

تعدّ عائلة \EN{GPT} من \EN{OpenAI} الأبرز في هذا المجال، حيث أطلقت الشركة
نموذج \EN{GPT-3.5} أواخر 2022 ثم \EN{GPT-4} في مارس 2023 \cite{openai2023gpt4}.
يُقدَّر أن \EN{GPT-4} يضم نحو 1.7 تريليون معامل، ويتميز بقدرات استدلالية
متقدمة وفهم سياقي عميق. وفي المنطقة العربية، أطلقت شركة \EN{Inception}
الإماراتية نموذج «\EN{Jais}» عام 2023 وهو أول نموذج لغوي عربي كبير، كما
أطلقت \EN{Alibaba} نموذج «\EN{Qwen}» الذي يدعم العربية بشكل جيد
\cite{elhaij2023arabic}.

تعتمد هذه النماذج على معمارية المحول (\EN{Transformer}) التي قدمتها ورقة
\EN{Vaswani} وآخرين عام 2017 \cite{vaswani2017attention}، والتي أحدثت ثورة من
خلال آلية الانتباه الذاتي (\EN{Self-Attention}). هذه الآلية تتيح للنموذج فهم
العلاقات بين الكلمات بغض النظر عن مواقعها في الجملة، مما يحسّن الفهم السياقي
بشكل كبير.

رغم قوتها الواضحة، تعاني النماذج اللغوية الكبيرة من قيود جوهرية: الهلوسة
(\EN{Hallucination}) أي توليد معلومات خاطئة بثقة عالية، وعدم القدرة على
الوصول للمعرفة الخاصة بالمجال، وقيود حجم السياق. هذه القيود هي ما دفع
الباحثين لتطوير تقنيات \EN{RAG} التي تُعالج لاحقاً في هذا الفصل.

\subsection{تقنية تحويل اللغة الطبيعية إلى \EN{SQL} (\EN{NL2SQL})}
\label{subsec:nl2sql}

تُعدّ تقنية \EN{NL2SQL} إحدى أهم تطبيقات الذكاء الاصطناعي في مجال قواعد
البيانات، حيث تُمكّن المستخدمين من الاستعلام عن البيانات بلغتهم الطبيعية دون
الحاجة لمعرفة \EN{SQL} \cite{yu2023spider}. وقد شهدت هذه التقنية قفزة نوعية
منذ ظهور النماذج اللغوية الكبيرة، حيث ارتفعت الدقة من نحو 35\% عام 2018 إلى
أكثر من 85\% عام 2024 على المعايير القياسية.

من أبرز الأنظمة الحالية في هذا المجال \EN{Vanna.ai}، وهو إطار عمل مفتوح
المصدر يدمج بين \EN{RAG} و \EN{NL2SQL}، ويُحقق دقة تتجاوز 80\% عند توفير
السياق المناسب \cite{vanna2024}. يعتمد \EN{Vanna} على تدريب نموذج خاص بكل
قاعدة بيانات، ويُستخدم على نطاق واسع في التطبيقات التجارية. كما أطلقت
\EN{Snowflake} أوائل 2025 نموذج \EN{Arctic-Text2SQL-R1} الذي حقق 71.83\% دقة
على معيار \EN{BIRD} الأكثر تحدياً، باستخدام تقنية التعلم المعزز في التدريب
\cite{snowflake2025arctic}.

من أهم التحديات في \EN{NL2SQL} ربط المخطط (\EN{Schema Linking})، حيث يجب
على النظام تحديد الجداول والأعمدة ذات الصلة بالسؤال من بين مئات العناصر في
قاعدة البيانات. ولهذا السبب، يلجأ كثير من المطورين إلى تبسيط المخطط عبر طبقات
وسيطة، وهو ما يبرر ابتكار طبقة \EN{Views} الذكية في مشروع بيان.

\subsection{تقنية الاسترجاع المعزز بالتوليد (\EN{RAG})}
\label{subsec:rag}

تُمثل تقنية \EN{RAG} إطاراً هجيناً يجمع بين قدرات النماذج اللغوية على التوليد،
وقواعد المعرفة الخارجية على توفير المعلومات الدقيقة \cite{lewis2020rag}. وقد
أصبحت هذه التقنية المعيار الذهبي لتطبيقات \EN{LLM} في البيئات المؤسسية منذ
عام 2023، حيث تُعالج مشكلة الهلوسة وتُمكّن النموذج من الوصول للمعرفة الخاصة
بالمجال.

يعمل نظام \EN{RAG} وفق سير عمل من ثلاث مراحل: الاسترجاع (\EN{Retrieval}) حيث
يتم تحويل سؤال المستخدم إلى متجه وبحث عن المستندات الأكثر صلة، ثم التوليد
المعزز (\EN{Augmented Generation}) حيث يُمرَّر المستندات المسترجعة مع السؤال
إلى النموذج لتوليد إجابة دقيقة، وأخيراً الاستشهاد (\EN{Citation}) حيث يُربط
كل جزء من الإجابة بالمصدر المستخرج منه.

\begin{figure}[!htbp]
  \centering
  \includegraphics[width=0.85\textwidth]{rag-flow.png}
  \caption{مخطط تدفق تقنية \EN{RAG} في نظام بيان.}
  \label{fig:rag-flow}
\end{figure}

من أبرز الأطر الحالية لبناء أنظمة \EN{RAG}، يأتي \EN{LangChain}
\cite{langchain2023} في المقدمة كإطار العمل الأكثر انتشاراً، ويقدم مكونات
قابلة للتخصيص للبحث والتوليد. ويأتي \EN{LlamaIndex} \cite{llamaindex2023}
كإطار متخصص في \EN{RAG}، يقدم أدوات متقدمة لفهرسة المستندات وإدارة السياق.
كما تقدم \EN{Microsoft} خدمة \EN{Azure AI Search} المتكاملة مع \EN{RAG}،
وتقدم \EN{Google} خدمة \EN{Vertex AI Search}.

يعتمد أداء \EN{RAG} بشكل جوهري على جودة نظام الاسترجاع، والذي يتطلب قاعدة
بيانات شعاعية (\EN{Vector Database}) متخصصة. وقد ظهرت عدة حلول مفتوحة المصدر
وتجارية في السنوات الأخيرة، أبرزها \EN{Qdrant} \cite{qdrant2024} الذي اعتمد
عليه مشروع بيان، و \EN{Pinecone} التجاري، و \EN{Milvus}، إضافة إلى امتداد
\EN{pgvector} لـ \EN{PostgreSQL}.

\subsection{الأنظمة متعددة الوكلاء (\EN{Multi-Agent Systems})}
\label{subsec:multi-agent}

تُعرَّف الأنظمة متعددة الوكلاء في سياق الذكاء الاصطناعي الحديث بأنها أنظمة
تُكلف فيها عدة وكلاء مستقلين بإجراء مهام متخصصة، مع وجود آلية تنسيق تضمن
تكامل مخرجاتهم. وقد ظهر هذا المفهوم بقوة منذ عام 2023 مع تطور قدرات النماذج
اللغوية، حيث أصبح كل وكيل عبارة عن \EN{LLM} مُوجَّه بمهمة وشخصية محددة.

يُعدّ إطار \EN{LangGraph} \cite{langgraph2024} من \EN{LangChain} الأبرز في
هذا المجال، ويستخدم نظرية المخططات لتمثيل تدفق العمل بين الوكلاء كعقد
(\EN{Nodes}) وحواف (\EN{Edges}). يتميز بتحكم دقيق في التدفق ودعم الحلقات
التكرارية والتفرعات المشروطة، وهو ما اعتمد عليه مشروع بيان. ويأتي إطار
\EN{AutoGen} \cite{wu2023autogen} من \EN{Microsoft Research} كخيار مرن قائم
على المحادثات بين الوكلاء، لكنه يعاني من زيادة في استهلاك الـ \EN{Tokens}.

إطار \EN{CrewAI} \cite{crewai2024} هو إطار حديث أطلق عام 2024، يركز على
محاكاة فرق العمل البشرية بأدوار محددة، ويتميز بسهولة الاستخدام. أما إطار
\EN{MetaGPT} \cite{hong2023metagpt} فيقدم نهجاً مختلفاً، حيث يحاكي فريق تطوير
برمجيات تقليدي (\EN{Product Manager, Architect, Engineer}) لإنتاج شفرات
برمجية متكاملة، وقد حقق تفوقاً ملحوظاً في إنتاج الكود الصحيح.

تُظهر هذه المراجعة أن \EN{LangGraph} هو الإطار الأنسب لمشروع بيان، نظراً
لكونه يدعم بناء أنظمة معقدة ذات وكلاء متخصصين، مع إمكانية تتبع مسار التنفيذ
(\EN{Agent Trace}) المطلوبة للشفافية في البيئة الجامعية.

\begin{figure}[!htbp]
  \centering
  \includegraphics[width=0.7\textwidth]{multi-agent-arch.png}
  \caption{معمارية النظام متعدد الوكلاء القائم على \EN{LangGraph}.}
  \label{fig:multi-agent}
\end{figure}

\subsection{تحديات معالجة اللغة العربية (\EN{Arabic NLP})}
\label{subsec:arabic-nlp}

تواجه معالجة اللغة العربية طبيعياً تحديات جوهرية تختلف عن نظيراتها في اللغة
الإنجليزية، وتعود هذه التحديات إلى عوامل لغوية وتقنية \cite{elhaij2023arabic}.
فمن الناحية اللغوية، تتميز العربية بالصرف الغني حيث يمكن اشتقاق عشرات الكلمات
من جذر واحد ثلاثي، مما يزيد من حجم المفردات وصعوبة التطابق. كما تتميز بتركيب
الجملة المرن الذي يصعّب التحليل النحوي الآلي، وظاهرة الإعراب التي تُغير الحركة
النهائية للكلمات.

أما من الناحية التقنية، فتعاني العربية من فجوة الموارد الحادة؛ إذ تشير
الإحصائيات إلى أن نسبة المحتوى العربي على الإنترنت لا تتجاوز 1.2\% من
الإجمالي، في حين أن المتحدثين بالعربية يشكلون نحو 5\% من سكان العالم
\cite{elhaij2023arabic}. وتعتمد معظم التطبيقات التجارية على نماذج إنجليزية
تُظهر أداءً ضعيفاً على العربية، وتشمل النماذج المدعومة للعربية: \EN{Jais}
النموذج الإماراتي المتخصص في العربية، و \EN{AraBERT} النموذج المبني على
\EN{BERT} والمتخصص في النصوص العربية، و \EN{Qwen 2.5} من \EN{Alibaba} الذي
يدعم العربية جزئياً ويتميز بكفاءة عالية في المهام البرمجية.

تتفاقم هذه التحديات في السياق اليمني تحديداً، حيث تُستخدم لهجة محلية تختلف عن
الفصحى والعاميات الخليجية والمصرية، وغير مدعومة بأي موارد برمجية معروفة.
ولهذا السبب، يعتمد مشروع بيان استراتيجية لغوية مرنة تقبل الفصحى والعامية، مع
الاستفادة من قدرة النماذج اللغوية الحديثة على التعميم بدلاً من تدريب نموذج
متخصص.

%------------------------------------------------------------------------------
\section{الأنظمة والدراسات السابقة}
\label{sec:related-work}

تم تصنيف الأنظمة والدراسات السابقة المراجعة في خمس فئات رئيسية، وفقاً لمجال
التخصص وارتباطها بمشروع بيان. وقد رُوعي في الاختيار التركيز على الأنظمة
الحديثة (2021-2025) والتجارية المتاحة فعلياً.

\subsection{الأنظمة في مجال تحويل اللغة الطبيعية إلى \EN{SQL}}
يُعدّ معيار \EN{Spider} الذي أُطلق عام 2018 المرجع القياسي لتقييم أنظمة
\EN{NL2SQL}، حيث يضم 10,181 سؤالاً عبر 200 قاعدة بيانات. وفي عام 2023،
أُطلق معيار \EN{BIRD} الأكثر تحدياً لتركيزه على قواعد البيانات الواقعية مع
قيم شاذة وأخطاء. حققت أفضل الأنظمة 40\% فقط على \EN{BIRD} مقابل 60\% على
\EN{Spider}، مما يُبرز فجوة الأداء في البيئات الواقعية \cite{yu2023spider}.

من أبرز الأنظمة الحالية في هذا المجال، يُبرز نظام \EN{Vanna.ai}
\cite{vanna2024} كأحد أكثر الأطر مفتوحة المصدر انتشاراً. يعتمد \EN{Vanna} على
دمج \EN{RAG} مع \EN{NL2SQL}، ويتدرب على قاعدة بيانات كل عميل لرفع الدقة.
أما نظام \EN{Snowflake Arctic-Text2SQL-R1} \cite{snowflake2025arctic} فهو
نموذج متخصص أطلق أوائل 2025، يستخدم تقنية التعلم المعزز في التدريب ويحقق
71.83\% على معيار \EN{BIRD}.

تجدر الإشارة إلى أن جميع هذه الأنظمة مصممة للبيئة الإنجليزية، ولا تقدم دعماً
كافياً للغة العربية. كما أنها تعمل كأنظمة أحادية (\EN{Single-Agent})، ولا
تستفيد من البنية متعددة الوكلاء في المهام المعقدة.

\subsection{الأنظمة في مجال الاسترجاع المعزز بالتوليد (\EN{RAG})}
منذ تقديم مفهوم \EN{RAG} بشكلها الحديث عام 2020 \cite{lewis2020rag}، شهدت
هذه التقنية تطوراً متسارعاً وأصبحت المعيار الذهبي لتطبيقات \EN{LLM} المؤسسية.
أظهرت الدراسات أن \EN{RAG} يقلل الهلوسة بنسبة تصل إلى 50\% مقارنة بالنماذج
اللغوية التقليدية، مما جعله خياراً مفضلاً للتطبيقات التي تتطلب دقة عالية.

من أبرز الأطر الحالية لبناء أنظمة \EN{RAG}، يأتي \EN{LangChain}
\cite{langchain2023} في المقدمة كإطار العمل الأكثر انتشاراً، ويقدم مكونات
قابلة للتخصيص للبحث والتوليد. ويأتي \EN{LlamaIndex} \cite{llamaindex2023}
كإطار متخصص في \EN{RAG}، يقدم أدوات متقدمة لفهرسة المستندات وإدارة السياق.
كما تقدم \EN{Microsoft} خدمة \EN{Azure AI Search} المتكاملة مع \EN{RAG}،
وتقدم \EN{Google} خدمة \EN{Vertex AI Search}.

على صعيد قواعد البيانات الشعاعية، يبرز \EN{Qdrant} \cite{qdrant2024} كحل
مفتوح المصدر يتميز بأداء عالٍ ودعم فلترة المتجهات بالبيانات الوصفية، وهو ما
اعتمد عليه مشروع بيان. كما تقدم \EN{Pinecone} حلاً تجارياً سحابياً، ويقدم
\EN{Milvus} حلاً مفتوحاً موزعاً. وفي سياق قواعد البيانات العلائقية، يقدم
امتداد \EN{pgvector} لـ \EN{PostgreSQL} إمكانية البحث الشعاعي دون الحاجة
لقاعدة بيانات منفصلة.

\subsection{الأنظمة في مجال الأنظمة متعددة الوكلاء}
رغم التطور المتسارع الذي شهدته الأنظمة متعددة الوكلاء منذ عام 2023 (انظر
المراجعة التفصيلية في القسم~\ref{subsec:multi-agent})، فإن التطبيقات المؤسسية
الجاهزة التي تعتمد هذه البنية لا تزال محدودة. تتركز الأطر المتاحة
(\EN{LangGraph}، \EN{AutoGen}، \EN{CrewAI}، \EN{MetaGPT}) على المستوى
التقني لبناء الوكلاء، دون تقديم حلول جاهزة للمؤسسات. هذا يعني أن أي مؤسسة
ترغب في اعتماد البنية متعددة الوكلاء مطالبة ببناء نظامها من الصفر، وهو ما يفسر
ندرة الأنظمة المؤسسية المبنية على هذه التقنية، ويبرر أهمية مشروع بيان الذي
يقدم تطبيقاً ميدانياً متكاملاً للبنية متعددة الوكلاء في بيئة جامعية.

\subsection{المساعدات الذكية في التعليم العالي}
رغم الانتشار الواسع للمساعدات الذكية في مجال خدمة العملاء، تبقى التطبيقات
المتخصصة في التعليم العالي محدودة. أطلقت \EN{OpenAI} عام 2024 خدمة
\EN{ChatGPT Edu} المخصصة للجامعات، التي تتيح بناء مساعدات مخصصة لكل مؤسسة.
ومع ذلك، تظل هذه الخدمة عامة ولا تدمج \EN{NL2SQL} للاستعلام عن قواعد البيانات
الجامعية.

تقدم \EN{Microsoft} خدمة \EN{Copilot for Education} مع تكاملات محدودة مع
أنظمة إدارة التعلم (\EN{LMS}) مثل \EN{Canvas} و \EN{Blackboard}، لكنها لا
تدعم العربية بشكل كامل. كما تقدم \EN{Anthropic} خدمة \EN{Claude for
Education} التي تركز على المساعدة الأكاديمية للطلاب.

على الصعيد العربي، أطلقت العديد من الجامعات مساعدات ذكية بسيطة للأسئلة
الشائعة، لكنها تعتمد على قواعد معرفة محدودة ولا تستفيد من تقنيات \EN{RAG}
أو \EN{NL2SQL}. وأكدت دراسة \EN{Al-Khalifa} وآخرين \cite{alkhalifa2023} أن
معظم الجامعات العربية تستخدم \EN{chatbots} بسيطة دون الاستفادة من التقنيات
المتقدمة.

تُبرز هذه المراجعة أن الميدان يفتقر إلى نظام متكامل يجمع بين \EN{NL2SQL} و
\EN{RAG} والأنظمة متعددة الوكلاء في بيئة جامعية عربية، وهو ما يسعى مشروع
بيان لسده.

\subsection{الدراسات العربية}
على الصعيد العربي، تظل الدراسات في مجال المساعدات الذكية المؤسسية محدودة
وتتسم بالتشتت بين أربع اتجاهات رئيسية لا تلتقي عادةً في بحث واحد: معالجة
اللغة العربية، استرجاع المعلومات العربية، \EN{NL2SQL} للعربية، وتطبيقات
الذكاء الاصطناعي في التعليم.

دراسة \EN{Magdy} وآخرين \cite{magdy2024arabicir} استعرضت واقع البحث في
استرجاع المعلومات العربية، وأكدت أن الأنظمة التجارية تعاني من ضعف 30-40\% في
الدقة على العربية مقارنة بالإنجليزية. كما أكدت دراسة \EN{El-Haj}
\cite{elhaij2023arabic} المرجعية في معالجة اللغة العربية طبيعياً، أن المشاريع
العربية تواجه ثلاث تحديات رئيسية: ندرة \EN{corpora} عربية عالية الجودة، عدم
وجود معايير قياسية موحدة، وضعف التمويل البحثي.

أما في مجال \EN{NL2SQL} العربية تحديداً، فالدراسات شبه منعدمة. أغلب الأبحاث
العربية في قواعد البيانات تركز على واجهات التعبير النمطي أو الترجمة الآلية
للاستعلامات، دون الاستفادة من النماذج اللغوية الكبيرة. وقد رصدت دراسة
\EN{Al-Khalifa} \cite{alkhalifa2023} غياب أي نظام \EN{NL2SQL} عربي منشور حتى
عام 2023، مما يجعل مشروع بيان من الرواد في هذا المجال.

في مجال المساعدات الذكية التعليمية العربية، أطلقت جامعة الملك سعود عام 2022
مساعداً ذكياً للأسئلة الأكاديمية اعتمد على قاعدة معرفة ثابتة دون \EN{RAG}،
كما طوّرت جامعة الملك عبدالعزيز نظاماً مشابهاً لإرشاد الطلاب. لكن هذه
المحاولات تظل بسيطة ولا تدمج تقنيات الاسترجاع المعزز بالتوليد، ولا تدعم
الاستعلام عن قواعد البيانات الجامعية.

في البيئة اليمنية تحديداً، تكاد تنعدم الدراسات المنشورة في مجال الذكاء
الاصطناعي المؤسسي. وأكد تقرير الأمم المتحدة للتحول الرقمي
\cite{un2024egd} أن اليمن تحتل المرتبة 175 عالمياً في مؤشر تطور الحكومة
الإلكترونية. ويُعزى ذلك إلى عوامل متعددة، تشمل: محدودية البنية التحتية
الرقمية، وضعف الاستثمار في البحث والتطوير، وندرة الكوادر المتخصصة. هذا الواقع
يُبرز أهمية مشروع بيان كمساهمة بحثية وتطبيقية في سد هذه الفجوة، حيث يقدم
أول نظام \EN{NL2SQL}+\EN{RAG} عربي مفتوح المصدر مخصص للبيئة الجامعية
اليمنية، مع توثيق كامل لمنهجيته وتقييم أدائه.

%------------------------------------------------------------------------------
\section{مقارنة الأنظمة السابقة}
\label{sec:comparison}

يقدم هذا القسم تحليلاً مقارناً للأنظمة والدراسات السابقة المراجعة، مع التركيز
على ستة أبعاد رئيسية: التقنيات المستخدمة، اللغة المدعومة، البيئة التطبيقية،
دقة \EN{NL2SQL}، دقة \EN{RAG}، والبنية المعمارية. ويُلخص
الجدول~\ref{tab:comparison} أبرز الأنظمة ومقارنتها بمشروع بيان.

\begin{table}[!htbp]
  \centering
  \caption{مقارنة شاملة بين الأنظمة السابقة ومشروع بيان.}
  \label{tab:comparison}
  \small
  \renewcommand{\arraystretch}{1.4}
  \begin{tabularx}{\textwidth}{|>{\raggedright\arraybackslash}p{2.5cm}|>{\raggedright\arraybackslash}X|>{\raggedright\arraybackslash}p{1.8cm}|>{\raggedright\arraybackslash}p{1.8cm}|>{\raggedright\arraybackslash}p{1.8cm}|>{\raggedright\arraybackslash}p{1.6cm}|>{\raggedright\arraybackslash}p{1.8cm}|}
    \hline
    \textbf{النظام/الإطار} & \textbf{التقنيات المستخدمة} & \textbf{اللغة} & \textbf{البيئة التطبيقية} & \textbf{دقة \EN{NL2SQL}} & \textbf{دقة \EN{RAG}} & \textbf{البنية} \\
    \hline
    \EN{Vanna.ai} \cite{vanna2024} (2024) & \EN{NL2SQL + RAG} & إنجليزية & تجارية & 80\%+ & غير مذكورة & أحادية \\
    \hline
    \EN{Arctic-Text2SQL} \cite{snowflake2025arctic} (2025) & \EN{NL2SQL + RL} & إنجليزية & تجارية & 71.83\% (\EN{BIRD}) & --- & أحادية \\
    \hline
    \EN{LangChain} \cite{langchain2023} (2023) & \EN{RAG + LLM} & متعددة & عامة & --- & عالية & أحادية \\
    \hline
    \EN{LlamaIndex} \cite{llamaindex2023} (2023) & \EN{RAG + LLM} & متعددة & عامة & --- & عالية & أحادية \\
    \hline
    \EN{AutoGen} \cite{wu2023autogen} (2023) & \EN{Multi-Agent} & إنجليزية & عامة & --- & --- & متعددة \\
    \hline
    \EN{CrewAI} \cite{crewai2024} (2024) & \EN{Multi-Agent} & إنجليزية & عامة & --- & --- & متعددة \\
    \hline
    \EN{MetaGPT} \cite{hong2023metagpt} (2024) & \EN{Multi-Agent} & إنجليزية & تطوير برمجي & --- & --- & متعددة \\
    \hline
    \EN{ChatGPT Edu} (2024) & \EN{LLM + RAG} & متعددة & تعليم عالٍ & --- & جيدة & أحادية \\
    \hline
    \EN{Qdrant} \cite{qdrant2024} (2024) & \EN{Vector DB} & متعددة & عامة & --- & عالية & --- \\
    \hline
    \textbf{مشروع بيان (المقترح)} & \EN{NL2SQL + RAG + Multi-Agent} & عربية (فصحى وعامية) & جامعي (يمني) & $\geq$ 85\% (مستهدف) & $\geq$ 80\% (مستهدف) & متعددة (\EN{LangGraph}) \\
    \hline
  \end{tabularx}
\end{table}

تكشف المقارنة عن عدة ملاحظات جوهرية. أولاً، تركز معظم الأنظمة في \EN{NL2SQL}
على البيئة الإنجليزية، مع غياب شبه تام للأنظمة العربية المتخصصة. ثانياً،
تعتمد الأنظمة المتقدمة مثل \EN{Vanna.ai} على دمج \EN{NL2SQL} و \EN{RAG}،
لكنها لا تستخدم البنية متعددة الوكلاء، مما يُحدها في المهام المعقدة. ثالثاً،
تفتقر التطبيقات في التعليم العالي إلى التكامل بين الاستعلام عن البيانات
الديناميكية والبحث في المستندات الثابتة.

كما يُلاحظ أن مشروع بيان هو الوحيد الذي يجمع بين التقنيات الأربع
(\EN{NL2SQL}، \EN{RAG}، \EN{Multi-Agent}، \EN{Arabic NLP}) في إطار موحد، مع
تخصيصه للبيئة الجامعية اليمنية. هذا التكامل يمثل نقطة قوة، لكنه في الوقت ذاته
يفرض تحديات تصميمية معقدة، خاصة في تنسيق الوكلاء وضمان جودة الأداء عبر
المكونات المختلفة.

%------------------------------------------------------------------------------
\section{الفجوة البحثية}
\label{sec:gap}

يتميز مشروع «بيان» بكونه النظام الوحيد — في حدود ما اطلعت عليه الدراسة — الذي
يجمع بين التقنيات الأربع الأساسية (\EN{NL2SQL}، \EN{RAG}، \EN{Multi-Agent
Systems}، \EN{Arabic NLP}) في إطار موحد ومحلي الاستضافة، مع تخصيصه المباشر
للبيئة المؤسسية والتعليمية. ورغم أن هذا التكامل يمثل نقطة قوة واضحة، إلا أنه
يفرض في الوقت ذاته تحديات تصميمية وتنفيذية معقدة، خاصة في تنسيق عمل الوكلاء
البرمجيين وضمان جودة الأداء والاعتمادية عبر المكونات المختلفة.

بناءً على المراجعة الشاملة للأنظمة والدراسات السابقة والتحليل المقارن، يمكن
تحديد الفجوة البحثية والتطبيقية التي يسعى مشروع «بيان» لسدها في \textbf{خمسة
محاور رئيسية}:

\textbf{المحور الأول — الفجوة التقنية:} رغم وجود أنظمة \EN{NL2SQL} متقدمة
(\EN{Vanna.ai}، \EN{Arctic-Text2SQL}) وأنظمة \EN{RAG} فعالة
(\EN{LangChain}، \EN{LlamaIndex})، إلا أنه لا يوجد نظام يجمع بين التقنيتين في
إطار موحد مع بنية متعددة الوكلاء. تعمل هذه الأنظمة بشكل منفرد، مما يُفقدها
التكامل المطلوب في البيئات المؤسسية المعقدة. مشروع بيان يسد هذه الفجوة عبر
الدمج العضوي للتقنيات الأربع (\EN{NL2SQL}، \EN{RAG}، \EN{Multi-Agent}،
\EN{Vector DB}).

\textbf{المحور الثاني — الفجوة اللغوية:} أكدت دراسات \EN{El-Haj}
\cite{elhaij2023arabic} و \EN{Magdy} \cite{magdy2024arabicir} على الفجوة
الواضحة في دعم اللغة العربية في أنظمة الذكاء الاصطناعي. تعاني الأنظمة
التجارية من انخفاض الدقة 30-40\% على العربية مقارنة بالإنجليزية، كما أنها غير
مصممة لفهم اللهجات المحلية. مشروع بيان يعالج هذه الفجوة عبر دعم العربية
الفصحى والعامية في واجهة محادثة \EN{RTL}، مع الاستفادة من نماذج \EN{Qwen} و
\EN{Jais} التي تدعم العربية.

\textbf{المحور الثالث — الفجوة التطبيقية} (\EN{Applied Gap}): تفتقر الأنظمة
التطبيقية في المؤسسات والقطاع الأكاديمي إلى التكامل بين الاستعلام عن البيانات
الرقمية والمالية (\EN{NL2SQL}) والبحث في اللوائح والسياسات والمستندات
التنظيمية (\EN{RAG}). تركز الأنظمة الحالية (مثل \EN{ChatGPT Enterprise} أو
\EN{ChatGPT Edu}) على جانب واحد فقط، مما يحد من قدرتها على تلبية الاحتياجات
التشغيلية الكاملة للمستفيدين. يقدم نظام «بيان» حلاً موحداً يدمج الجانبين عبر
وكلاء متخصصين.

\textbf{المحور الرابع — الفجوة السياقية:} الأنظمة التجارية مثل \EN{Vanna.ai}
و \EN{ChatGPT Edu} مصممة لبيئات تجارية غربية، ولا تأخذ في الاعتبار خصوصية
البيئة الجامعية اليمنية (اللوائح المحلية، اللهجة، البنية التحتية المحدودة).
مشروع بيان مصمم خصيصاً لهذه البيئة، مع استخدام \EN{Odoo ERP} كبيئة محاكاة
واقعية تُحقن ببيانات تعادل 15 عاماً.

\textbf{المحور الخامس — فجوة الابتكار:} لا توجد دراسة سابقة قدمت طبقة
\EN{Views} ذكية كآلية لتبسيط مخطط قاعدة البيانات وتحسين دقة \EN{NL2SQL}.
هذا الابتكار يمثل مساهمة بحثية أصلية يمكن البناء عليها في دراسات مستقبلية،
وقد يفتح آفاقاً جديدة في مجال \EN{Database Abstraction Layer}.

يُلخص الجدول~\ref{tab:gap} الفجوة البحثية ومساهمة مشروع بيان في سدها عبر
المحاور الخمسة.

\begin{table}[!htbp]
  \centering
  \caption{ملخص الفجوة البحثية ومساهمة مشروع بيان في سدها.}
  \label{tab:gap}
  \renewcommand{\arraystretch}{1.5}
  \begin{tabularx}{\textwidth}{|>{\raggedright\arraybackslash}p{2.5cm}|>{\raggedright\arraybackslash}X|>{\raggedright\arraybackslash}X|}
    \hline
    \textbf{المحور} & \textbf{الفجوة المحددة} & \textbf{مساهمة مشروع بيان} \\
    \hline
    تقنية & غياب الدمج بين \EN{NL2SQL} و \EN{RAG} و \EN{Multi-Agent} & دمج التقنيات الأربع في إطار موحد \\
    \hline
    لغوية & ضعف دعم العربية واللهجات & دعم فصحى وعامية في واجهة \EN{RTL} \\
    \hline
    تطبيقية & فصل الاستعلام عن البيانات عن البحث في المستندات & نظام موحد يجمع الجانبين \\
    \hline
    سياقية & غياب حلول مخصصة للبيئة المؤسسية والمحلية & تخصيص كامل للبيئة المؤسسية مع بيئة محاكاة لأكثر من نوع من قواعد البيانات \\
    \hline
    ابتكارية & غياب طبقة تجريدية مبسطة فوق قواعد البيانات & ابتكار طبقة \EN{Views} ذكية \\
    \hline
  \end{tabularx}
\end{table}

%------------------------------------------------------------------------------
\section{الخلاصة}
\label{sec:summary}

استعرض هذا الفصل الإطار النظري الشامل الذي يستند إليه \textbf{«مشروع المساعد
الذكي المؤسسي (بيان)»}، بدءاً من المفاهيم الأساسية في النماذج اللغوية الكبيرة
(\EN{LLMs})، مروراً بالتقنيات المتخصصة في \EN{NL2SQL} و \EN{RAG} والأنظمة
متعددة الوكلاء (\EN{Multi-Agent Systems})، وصولاً إلى تحديات معالجة اللغة
العربية طبيعياً (\EN{Arabic NLP}). وقد كشف الاستعراض النظري أن المشروع يعتمد
على ركائز علمية متينة ومثبتة، مع الاستفادة من أحدث الأطر والأدوات المتاحة
(2021-2025) في كل مجال.

كما قدم الفصل مراجعة شاملة للأنظمة والدراسات السابقة ومقارنة مرجعية بينها.
وأظهرت المراجعة أن مجال \EN{NL2SQL} شهد تقدماً كبيراً في السنوات الثلاث
الأخيرة مع ظهور أطر مثل \EN{Vanna.ai} و \EN{Arctic-Text2SQL}. وفي مجال
\EN{RAG}، أثبتت الأطر المتقدمة مثل \EN{LangChain} و \EN{LlamaIndex} كفاءتها
العالية. وفي مجال الأنظمة متعددة الوكلاء، أظهر إطار \textbf{\EN{LangGraph}}
تفوقاً ملحوظاً في إدارة الحالات (\EN{State Management}) والتحكم في مسارات
التنفيذ.

كشف التحليل المقارن أن مشروع «بيان» يمثل خطوة متقدمة في الدمج الشامل للتقنيات
الأربع في بيئة عمل مؤسسية، حيث اتُخذت جامعة العلوم والتكنولوجيا كدراسة حالة
وتطبيق تجريبي. وخلص الفصل إلى تحديد خمسة محاور رئيسية للفجوة البحثية (تقنية،
لغوية، تطبيقية، سياقية، وابتكارية) يسعى المشروع لسدها، مما يُبرر القرارات
التصميمية والهندسية المتخذة، ويهيئ القارئ للفصل الثالث الذي يقدم تحليلاً
تفصيلياً للمتطلبات والمنهجية المتبعة في تطوير النظام.
